\documentclass[11pt,a4paper]{article}
\usepackage[margin=0.85in,top=1in,bottom=1in]{geometry}
\usepackage{amsmath,amssymb,amsfonts}
\usepackage{graphicx}
\usepackage{float}
\usepackage{enumitem}
\usepackage{microtype}
\usepackage{caption}
\usepackage[hidelinks]{hyperref}
\usepackage[most,breakable]{tcolorbox}
\tcbuselibrary{breakable,skins}
\usepackage[utf8]{inputenc}
\usepackage[T1]{fontenc}
\usepackage{lmodern}
\captionsetup{font=small,labelfont=bf,skip=4pt}
\setlist{leftmargin=1.5em,topsep=2pt}

%% Breakable card — prevents content cutoff across pages
\newtcolorbox{solcard}[3]{
  breakable, enhanced,
  colback=white, colframe=gray!50, arc=2pt,
  title={\textbf{#1}\hfill\textsf{\small #3\,pts}},
  fonttitle=\normalsize\sffamily\bfseries,
  before upper={\small\textit{#2}\par\smallskip\hrule height 0.4pt\smallskip},
  top=4pt, bottom=6pt, left=7pt, right=7pt,
  parbox=false,
}

\title{\textbf{Roche limits --- Solution}}
\author{\normalsize Y22 (2022) — English translation}
\date{}
\begin{document}\maketitle\vspace{-0.5em}
\begin{solcard}{A1}{Find the value of the radius of the circular orbit $R_\text{BPS}$ at which the mechanical stress in the area of contact of the objects is equal to zero. The quantity $R_\text{BPS}$ is the Roche limit $R_{Roche}$ for a binary planetary system}{1.50}

Let $\omega$ be the angular velocity of rotation of the system. Then from the theorem on the motion of the center of mass we obtain: $$2m\omega^2R=GMm\left(\displaystyle\frac{1}{(R-r)^2}+\displaystyle\frac{1}{(R+r)^2}\right) $$ Consider separately the condition for the motion of the center of the near object on a circle of radius $R-r$: $$m\omega^2(R-r)=\displaystyle\frac{GMm}{(R-r)^2}+N-\displaystyle\frac{Gm^2}{4r^2} $$ Combining, we find: $$N=\displaystyle\frac{Gm^2}{4r^2}+\displaystyle\frac{GMm}{2R}\left(\displaystyle\frac{R-r}{(R+r)^2}-\displaystyle\frac{R+r}{(R-r)^2}\right) $$ Taking into account that $R_0\gg{r}$, find: $$N\approx{\displaystyle\frac{Gm^2}{4r^2}-\displaystyle\frac{3GMmr}{R^3}} $$ Force mechanical interaction between of the objects is non-negative for $$R\geq{r\sqrt[3]{\displaystyle\frac{12M}{m}}} $$ Then for $R_{DPS}$ we have:

\par\smallskip\noindent\textbf{Answer:} $$R_\text{BPS}=r\sqrt[3]{\displaystyle\frac{12M}{m}} $$
\end{solcard}

\begin{solcard}{B1}{Obtain an exact expression for the total potential $\varphi_\text{ext}$ in the field of the centrifugal force and the gravitational field of the planet. Express your answer in terms of $x$, $y$, $z$, $G$, $\rho_M$, $R_0$ and $R$.}{0.70}

The orbital angular velocity of the satellite $\omega$ is equal to $$\omega=\sqrt{\cfrac{GM}{R^3}} $$ where $M=\cfrac{4\pi R^3_0\rho_M}{3}$. Potential energy particle mass $m$ in the field of the centrifugal force (consider it equal to zero on the axis of rotation) equals $$W_p=-\frac{m\omega^2r^2_{\perp}}{2} $$ where $r_{\perp}$ - distance from particle to the axis of rotation. Then for potential have: $$\varphi=-\cfrac{4\pi R^3_0\rho_MG}{3R}\left(\cfrac{R}{r}+\cfrac{r^2_\perp}{2R^2}\right)-\varphi_O $$ from where:

\par\smallskip\noindent\textbf{Answer:} $$\varphi=\frac{2\pi R^3_0\rho_MG}{R}-\cfrac{4\pi R^3_0\rho_MG}{3R}\left(\cfrac{R}{\sqrt{(R+x)^2+y^2+z^2}}+\cfrac{(R+x)^2+y^2}{2R^2}\right) $$
\end{solcard}

\begin{solcard}{B2}{In the $x,y,z\ll{R}$ approximation, the expression for the $\varphi$ potential can be represented as follows: $$\varphi_{\text{ext}}=A_1x^2+B_1y^2+C_1z^2 $$ Find coefficient $A_1$, $B_1$ and $C_1$. Answer express through $R$, $G$, $R_0$ and $\rho_M$.}{0.80}

For the expansion of $\cfrac{R}{\sqrt{(R+x)^2+y^2+z^2}}$ we have: $$\cfrac{R}{\sqrt{(R+x)^2+y^2+z^2}}\approx{1-\frac{x}{R}-\frac{x^2+y^2+z^2}{2R^2}}+\frac{3x^2}{2R^2} $$ Substitute in expression for $\varphi$, obtain: $$\varphi=-\cfrac{2\pi R^3_0\rho_MG}{R^3}x^2+\cfrac{2\pi R^3_0\rho_MG}{3R^3}z^2 $$ Thus:

\par\smallskip\noindent\textbf{Answer:} $$A=-\cfrac{2\pi R^3_0\rho_MG}{R^3} \quad B=0 \quad C=\cfrac{2\pi R^3_0\rho_MG}{3R^3} $$
\end{solcard}

\begin{solcard}{C1}{Prove that the gravitational field inside the cavity is zero.}{0.80}

Let us turn to the proof of equality to zero in the cavity of the ball concentric to it. Let us consider some point of the cavity $A$. Let us select two elementary volumes $dV_1$ and $dV_2$, the fields of which compensate each other: $$\frac{dV_1}{r^2_1}=\frac{dV_2}{r^2_2} $$ An arbitrary ellipsoid can be obtained by compression from a sphere along mutually perpendicular axes. Since the volume of any elementary parallelepiped is $dV = dx\,dy\,dz$, all elementary volumes decrease by the same number of times. The ratios of distances between points will also be preserved. Then, since the total gravitational field inside the cavity of the ball is zero, the gravitational field inside the cavity of the ellipsoid is also equal to zero.

\end{solcard}

\begin{solcard}{C2}{Let the potential at the center of the ellipsoid be $\varphi_0$. Consider a point inside the ellipsoid, the radius vector relative to the center is $\vec{r}$. The potential at a given point is $\varphi_1$. What is the potential $\varphi_2=\varphi(\gamma\vec{r})$ at a point with radius vector $\gamma\vec{r}$ also located inside the ellipsoid? Express your answer using $\varphi_0$, $\varphi_1$ and $\gamma$.}{0.40}

Consider the points $1$ and $2$ lying on the same straight line passing through the center of the ellipsoid. Let the distances from the center of the ellipsoid to these points is equal to to $r_1$ and $r_2$, respectively, and the straight line under consideration makes an angle $\alpha$ with the axis $x$. Both points belong to the boundaries of ellipsoids similar to the one under consideration with some coefficients $k_1$ and $k_2$. Since the gravitational field inside ellipsoidal cavities is zero, the potential difference between the point $1$ and the center of the ellipsoid in question $\varphi_1-\varphi_0=f(\alpha)r^2_1$, as well as $\varphi_2-\varphi_0=f(\alpha)r^2_2$, where $f(\alpha)$ is a certain function depending only on the angle $\alpha$. From here we find:

\par\smallskip\noindent\textbf{Answer:} $$\varphi_2=\varphi_0+(\varphi_1-\varphi_0)\gamma^2 $$
\end{solcard}

\begin{solcard}{C3}{Show that at a point on an ellipsoid with coordinates $x,y,z$ the potential can be represented as: $$\varphi=\varphi_0+A_2x^2+B_2(y^2+z^2) $$ $\textit{Note:}$ note that this expression is applicable for ellipsoids of arbitrarily small sizes}{1.00}

The distances $r$ can be arbitrarily small. Let's consider points close to the center of the ellipsoid. For them: $$\vec{g}=\frac{\partial\vec{g}}{\partial x}dx+\frac{\partial\vec{g}}{\partial y}dy+\frac{\partial\vec{g}}{\partial z}dz $$ where the derivatives are taken at the origin of coordinates. Since ellipsoid - is a symmetric figure, for a point lying on the $x$, $y$ and $z$ axes correct correspondingly. $$\frac{\partial\vec{g}}{\partial x}=\frac{\partial g_x}{\partial x}\vec{e}_x\quad\frac{\partial\vec{g}}{\partial y}=\frac{\partial g_y}{\partial y}\vec{e}_y\quad\frac{\partial\vec{g}}{\partial z}=\frac{\partial g_z}{\partial z}\vec{e}_z $$ hence find: $$\frac{\partial g_r}{\partial r}=-2f(\alpha)=(\vec{e}_r\cdot{\vec{e}_x})\frac{\partial g_x}{\partial x}\frac{\partial x}{\partial r}+(\vec{e}_r\cdot{\vec{e}_y})\frac{\partial g_y}{\partial y}\frac{\partial y}{\partial r}+(\vec{e}_r\cdot{\vec{e}_z})\frac{\partial g_z}{\partial z}\frac{\partial z}{\partial r} $$ Let $g_\perp$ - component gravitational field ellipsoid, perpendicular axis $x$.


Considering that for an ellipsoid of revolution $\cfrac{\partial g_y}{\partial y}=\cfrac{\partial g_z}{\partial z}=\cfrac{\partial g_\perp}{\partial r_\perp}$, we obtain $$-2f(\alpha)=\frac{\partial g_x}{\partial x}\cos^2\alpha+\frac{\partial g_\perp}{\partial r_\perp}\sin^2\alpha $$ Since relation correct for any $\alpha$, all partial derivatives must remain constant at any point of the ellipsoid. Substituting into the expression for the potential, we find:

\par\smallskip\noindent\textbf{Answer:} $$\varphi=\varphi_0-\frac{\partial g_x}{\partial x}\frac{x^2}{2}-\frac{\partial g_\perp}{\partial r_\perp}\frac{(y^2+z^2)}{2} $$
\begin{figure}[H]
  \centering
  \includegraphics[width=0.72\textwidth]{T5_S_EN_assets/fig1.jpg}
\end{figure}
\end{solcard}

\begin{solcard}{C4}{Express $B_2$ in terms of $A_2$, $G$ and $\rho$. $\textit{Note:}~$ use Gauss's theorem in differential form for the gravitational field: $$\displaystyle\frac{\partial^2{\varphi}}{\partial{x}^2}+\displaystyle\frac{\partial^2{\varphi}}{\partial{y}^2}+\displaystyle\frac{\partial^2{\varphi}}{\partial{z}^2}=4\pi\rho G $$}{0.30}

From Gauss’s theorem in differential form we find: $$\frac{\partial g_x}{\partial x}+\frac{\partial g_y}{\partial y}+\frac{\partial g_z}{\partial z}=-2A_2-4B_2=-4\pi G\rho $$ from where:

\par\smallskip\noindent\textbf{Answer:} $$B_2=\pi G\rho-\frac{A_2}{2} $$
\end{solcard}

\begin{solcard}{C5}{The potential $\varphi$ on the surface of the ellipsoid can be represented as a function of only the coordinate $x$ in the following form: $$\varphi=C_2+x^2D_2 $$ Express $C_2$ and $D_2$ through $\varphi_0$, $A_2$, $G$, $\rho$, $a$, $x$ and eccentricity $e$.}{0.30}

$$y^2+z^2=a^2(1-e^2)-(1-e^2)x^2 $$ from which: $$\varphi=\varphi_0+\left(\pi\rho G-\cfrac{A_2}{2}\right)a^2(1-e^2)+\left(\cfrac{A_2(3-e^2)}{2}-\pi\rho G(1-e^2)\right)x^2 $$

\par\smallskip\noindent\textbf{Answer:} $$C_2=\varphi_0+\left(\pi\rho G-\cfrac{A_2}{2}\right)a^2(1-e^2) \quad D_2=\cfrac{A_2(3-e^2)}{2}-\pi\rho G(1-e^2) $$
\end{solcard}

\begin{solcard}{C6}{Let us solve an auxiliary problem: a mass with surface density $\sigma$ is uniformly distributed over the surface of a disk of radius $R$. Find the potential of the disk $\varphi_\text{further}$ on the axis at a distance $x$ from its center. Express your answer using $\sigma$, $G$, $R$ and $x$.}{0.50}

Let's split the disk into rings: $$\varphi_\text{further}=-\pi\sigma G\int\limits_0^R\frac{2rdr}{\sqrt{r^2+x^2}} $$

\par\smallskip\noindent\textbf{Answer:} $$\varphi_\text{further}=-2\pi\sigma G\left(\sqrt{R^2+x^2}-x\right) $$
\end{solcard}

\begin{solcard}{C7}{Find the potential $\varphi_0$ at the center of the ellipsoid. Express your answer in terms of $G$, $\rho$, $a$ and $e$. $\textit{Note:}$ you may need the following integral: $$\displaystyle\int\sqrt{1+x^2}dx=\cfrac{x\sqrt{1+x^2}+\ln\left(x+\sqrt{1+x^2}\right)}{2}+C $$}{1.00}

Let's divide the ellipsoid into disks, taking into account that $d\sigma=\rho dx$: $$\varphi_0=-2\pi\rho G\cdot{2\int\limits_{0}^a\left(\sqrt{x^2+y^2+z^2}-x\right)dx} $$ Integral from second term equal $-\cfrac{a^2}{2}$. Substitute $y^2+z^2$ from the equation ellipsoid: $$\varphi=2\pi\rho Ga^2-4\pi\rho G\int\limits_{-a}^a\sqrt{b^2+\frac{c^2x^2}{a^2}}dx=2\pi\rho Ga^2-\frac{4\pi\rho G ab^2}{c}\int\limits_0^{\frac{c}{b}}\sqrt{1+t^2}dt $$ Using the antiderivative: $$\varphi_0=2\pi\rho Ga^2-\frac{2\pi\rho G ab^2\left(\cfrac{ac}{b^2}+\ln\left(\cfrac{a+c}{b}\right)\right)}{c} $$ Let us express it through eccentricity:

\par\smallskip\noindent\textbf{Answer:} $$\varphi_0=-\frac{2\pi\rho Ga^2(1-e^2)}{e}\ln\sqrt{\cfrac{1+e}{1-e}} $$
\end{solcard}

\begin{solcard}{C8}{Find the potential of the ellipsoid $\varphi_a$ at the point with coordinates $(x, y, z)=(a, 0, 0)$. Express your answer using $G$, $\rho$, $a$ and $e$. $\textit{Note:}$ you may need the following integral: $$\displaystyle\int\sqrt{x^2-1}dx=\cfrac{x\sqrt{x^2-1}-\ln\left(x+\sqrt{x^2-1}\right)}{2}+C $$}{1.00}

Let's again divide the ellipsoid into disks: $$\varphi_a=-2\pi\rho G\int\limits_{-a}^{a}\left(\sqrt{(a+x)^2+y^2+z^2}-(x+a)\right)dx $$ Integral from second term equal $-2a^2$. Substitute $y^2+z^2$ from the equation ellipsoid and transform: $$\int\limits_{-a}^{a}\sqrt{(a+x)^2+y^2+z^2}dx=\int\limits_{-a}^{a}\sqrt{\left(\frac{cx}{a}+\frac{a^2}{c}\right)^2-\frac{b^4}{c^2}}dx=\frac{ab^4}{c^3}\int\limits_{\frac{a^2-c^2}{b^2}}^{\frac{a^2+c^2}{b^2}}\sqrt{t^2-1}dt= $$ $$=\frac{a^2(1-e^2)^2}{e^3}\int\limits_{1}^{\frac{1+e^2}{1-e^2}}\sqrt{t^2-1}dt=a^2\left(\frac{(1+e^2)}{e^2}-\frac{(1-e^2)^2}{e^3}\ln\sqrt{\frac{1+e}{1-e}}\right) $$ from where:

\par\smallskip\noindent\textbf{Answer:} $$\varphi_a=-2\pi\rho Ga^2\left(\frac{(1-e^2)}{e^2}-\frac{(1-e^2)^2}{e^3}\ln\sqrt{\frac{1+e}{1-e}}\right) $$
\end{solcard}

\begin{solcard}{C9}{Find the coefficient $A_2$. Express your answer in terms of $G$, $\rho$ and $e$.}{0.20}

$$\varphi_a-\varphi_0=A_2a^2=2\pi\rho Ga^2\left(\frac{1-e^2}{e^3}\ln\sqrt{\frac{1+e}{1-e}}-\frac{(1-e^2)}{e^2}\right) $$ from

\par\smallskip\noindent\textbf{Answer:} $$A_2=2\pi\rho G\left(\frac{1-e^2}{e^3}\ln\sqrt{\frac{1+e}{1-e}}-\frac{(1-e^2)}{e^2}\right) $$
\end{solcard}

\begin{solcard}{D1}{On a sheet of graph paper, plot the dependence of the coefficient $D_2$ on the eccentricity $e$ of the ellipsoid formed by the satellite.}{0.70}

Let's substitute $A_2$ into the expression for $D_2$: $$D_2=\pi\rho G~\frac{1-e^2}{e^3}\left((3-e^2)\ln\sqrt{\frac{1+e}{1-e}}-3e\right) $$

\begin{figure}[H]
  \centering
  \includegraphics[width=0.72\textwidth]{T5_S_EN_assets/fig2.jpg}
\end{figure}
\end{solcard}

\begin{solcard}{D2}{At what value of eccentricity $e$ is the Roche limit reached?}{0.30}

When the function $D_2$ reaches its maximum, a further decrease in the radius of the circular orbit is impossible, so the satellite is destroyed. The thought of solving an equation for a zero derivative is not reasonable.


From the graph we find:

\par\smallskip\noindent\textbf{Answer:} $$e=0,86\pm{0,02} $$
\begin{figure}[H]
  \centering
  \includegraphics[width=0.72\textwidth]{T5_S_EN_assets/fig3.jpg}
\end{figure}
\end{solcard}

\begin{solcard}{D3}{Find the Roche limit $R_{Roche}$. Express your answer in terms of $R_0$, $\rho_M$ and $\rho_m$.}{0.50}

The value of $D_2$ at the extremum is equal to $$D_{2_{max}}\approx{0,14\pi\rho_m G} $$ from where:

\par\smallskip\noindent\textbf{Answer:} $$R_{Roche}\approx{2,43R_0\sqrt[3]{\frac{\rho_M}{\rho_m}}} $$
\end{solcard}

\end{document}